% !TeX root = ../../Book4.tex
% 纲要标题：### 14.1 滤过与伴随分次
% 纲要位置：计划纲要.md，第 2827 行。

\section{滤过与伴随分次}
\label{b4:sec:1401}

滤过先把复形分成较简单的层，再保留微分在层间移动的信息。递增与递减编号会改变后面公式的方向，有界、穷尽与分离条件也必须分别说明。

% 纲要内容：
%
% * 递增、递减滤过
% * 有界、穷尽和分离条件
% * 首先固定指标方向
%

\subsection{递增滤过与递减滤过}
\label{b4:subsec:1401:01}

把一个复杂复形分成几层，往往比直接求核再求商容易。例如双复形的总次数固定以后，可以先保留横指标较大的那些项。微分可能把一层送到更深的一层；若暂时忽略这种移动，便得到一组较简单的复形。谱序列记录随后怎样把这些层间信息逐步补回来。

\begin{definition}\label{b4:def:filtered-complex}
设 \(C\) 为含幺环 \(R\) 上左模的上链复形。对每个整数 \(p\)，给定子复形 \(F^pC\subseteq C\)，若
\[
 F^{p+1}C\subseteq F^pC,
\]
则称 \((C,F)\) 为带\term[dijian-lvguo]{递减滤过}[decreasing filtration]的复形。其\term[bansui-fenci]{伴随分次}[associated graded object]为复形族
\[
 \operatorname{gr}_F^pC=F^pC/F^{p+1}C.
\]
若改用 \(F_pC\subseteq F_{p+1}C\)，则得到递增滤过。滤过映射 \(f:(C,F)\to(D,F)\) 是满足 \(f(F^pC)\subseteq F^pD\) 的复形映射。
\end{definition}

用 \(F^pC=F_{-p}C\) 可以互换递增、递减记号，但微分的指标也必须相应改变。本章固定递减约定。若 \(D^{a,b}\) 是第一象限双复形，令
\[
 F^p\operatorname{Tot}(D)^n=\bigoplus_{a\ge p}D^{a,n-a}.
\]
纵微分保持 \(a\)，横微分增大 \(a\)，所以总微分保持这个子对象。伴随分次只剩第 \(p\) 列，最先显现的是纵向上同调。

\subsection{有界、穷尽与分离条件}
\label{b4:subsec:1401:02}

\begin{definition}\label{b4:def:finite-filtration}
设 \(R\) 为环，\((C,F)\) 是左 \(R\)-模的递减滤过上链复形。若对每个总次数 \(n\) 都有 \(\bigcup_pF^pC^n=C^n\)，则称滤过穷尽；若对每个 \(n\) 都有 \(\bigcap_pF^pC^n=0\)，则称滤过分离。若对每个 \(n\) 存在整数 \(a_n<b_n\)，使 \(F^{a_n}C^n=C^n\)、\(F^{b_n}C^n=0\)，则称滤过逐度有限。若 \(a_n,b_n\) 可独立于 \(n\) 选取，则称滤过一致有界。
\end{definition}

逐度有限蕴含穷尽和分离，但不要求复形只有有限多个次数。第一象限双复形的列滤过在总次数 \(n\ge0\) 满足 \(F^0=C^n\)、\(F^{n+1}=0\)，就是本书最常用的逐度有限情形。

\ProseParagraphBreak
穷尽和分离本身不保证有限计算。例如 \(C^0=\mathbb Z\)、其他次数为零，取 \(F^pC^0=2^p\mathbb Z\)（\(p\ge0\)）、\(F^pC^0=\mathbb Z\)（\(p<0\)），便得到穷尽、分离但不逐度有限的滤过。商群 \(C^0/F^pC^0\) 的相容系统还包含并非来自某个整数的元素。因此讨论无限滤过时，是否需要完备化是实质问题；本章的主要收敛定理采用逐度有限条件，不把这些差别省去。

\subsection{滤过的指标方向}
\label{b4:subsec:1401:03}

总次数写作 \(n=p+q\)，其中 \(p\) 是滤过次数，\(q\) 只是补足总次数的第二指标。一个 \(x\in F^pC^{p+q}\) 若满足 \(\mathrm{d}x\in F^{p+r}C^{p+q+1}\)，便表示它的微分直到相隔 \(r\) 层才可能被看见。于是第 \(r\) 次出现的微分必须具有方向
\[
 \mathrm{d}_r:E_r^{p,q}\longrightarrow E_r^{p+r,q-r+1}.
\]
它使总次数增加一；在第零页竖直向上，在第一页水平向右，第二页向右两格、向下一格。

\begin{definition}\label{b4:def:spectral-sequence}
设 \(R\) 为含幺环。给定双分次左 \(R\)-模 \(E_r^{p,q}\)，其中 \(r\ge0\)、\(p,q\in\mathbb Z\)，以及 \(R\) 线性映射 \(\mathrm{d}_r:E_r^{p,q}\to E_r^{p+r,q-r+1}\)。若这些映射及指定同构满足
\[
 \mathrm{d}_r^2=0,\qquad E_{r+1}^{p,q}\simeq
 \frac{\ker(\mathrm{d}_r:E_r^{p,q}\to E_r^{p+r,q-r+1})}
 {\operatorname{im}(\mathrm{d}_r:E_r^{p-r,q+r-1}\to E_r^{p,q})}
\]
则称这些数据构成从第零页开始的上同调型\term[pu-xulie]{谱序列}[spectral sequence]。也可从某个较晚页开始给出这些数据。谱序列态射在各页保持双次数、与微分及指定同构相容。
\end{definition}
\symidx{\(E_r^{p,q},\mathrm{d}_r\)：谱序列的页与微分}

这个定义描述的是一列反复取上同调的对象，尚未指定它计算什么。目标及其滤过必须另外给出；不能只凭写出 \(E_2\) 就宣称已经算出某个上同调群。下面从实际滤过复形构造全部各页，再证明它们与原复形上同调的关系。一般交换范畴中的构造以子对象的交、像和商代替元素集合；证明中的有限提升可依\cref{b4:lem:generalized-element-calculus}解释，故逐度有限的结论也适用于交换范畴。
